\chapter{HASIL DAN PEMBAHASAN}
\label{ch:hasil_pembahasan}
\section{Parameter Fisis Fundamental: Studi Kasus Al dan Ca}
\label{sec:parameter_fisis}
Untuk memberikan gambaran transparan mengenai proses komputasi yang dilakukan, bagian ini menyajikan perhitungan langkah demi langkah untuk parameter fisis fundamental. Perhitungan ini menunjukkan bagaimana spektrum dihasilkan dari prinsip pertama, menggunakan data dari basis data NIST. Aluminium (Al) dan Kalsium (Ca) dipilih sebagai studi kasus representatif pada kondisi plasma T = 8000 K dan $n_e = 1.0 \times 10^{17}$ cm$^{-3}$.
\subsection{Keseimbangan Ionisasi Makroskopis (Analisis Saha)}

Langkah pertama adalah menentukan keadaan ionisasi makroskopis dari plasma menggunakan Persamaan Saha. Tabel~\ref{tab:saha_results} menunjukkan hasil perhitungan fraksi atom netral dan terionisasi. Pada kondisi ini, kedua elemen didominasi oleh keadaan terionisasi (Al II dan Ca II), dengan fraksi ion masing-masing sebesar 85.4\% dan 83.0\%. Fraksi ini akan bertindak sebagai faktor skala utama yang menentukan kontribusi total dari setiap keadaan ionisasi pada spektrum akhir.

\begin{table}[H]
    \centering
    \caption{Hasil perhitungan Persamaan Saha untuk Al dan Ca pada T=8000 K dan $n_e=1.0 \times 10^{17}$ cm$^{-3}$.}
    \label{tab:saha_results}
    \small
    \begin{tabular}{lcccc}
        \toprule
        \textbf{Elemen} & \textbf{Energi Ionisasi (eV)} & \textbf{Rasio Saha ($N_{II}/N_{I}$)} & \textbf{Fraksi Netral (\%)} & \textbf{Fraksi Ion (\%)} \\
        \midrule
        Al & 5.986 & $5.86 \times 10^0$ & 14.58 & 85.42 \\
        Ca & 6.113 & $4.87 \times 10^0$ & 17.04 & 82.96 \\
        \bottomrule
    \end{tabular}
\end{table}
\subsection{Parameter Mikroskopis dan Perhitungan Intensitas}

Langkah kedua adalah menganalisis parameter mikroskopis. Tabel~\ref{tab:lines_and_partition} menyajikan nilai fungsi partisi (Z) yang dihitung untuk setiap keadaan ion, bersama dengan parameter atom untuk lima garis transisi terkuat yang dipilih berdasarkan intensitas relatif finalnya.

\begin{table}[H]
    \centering
    \caption{Fungsi partisi (Z) dan parameter transisi untuk 5 garis terkuat dari setiap keadaan ion Al dan Ca.}
    \label{tab:lines_and_partition}
    \small
    \begin{tabular}{lcccccc}
        \toprule
        \textbf{Elemen} & \textbf{$\lambda$ (nm)} & \textbf{$E_i$ (eV)} & \textbf{$E_k$ (eV)} & \textbf{$g_i$} & \textbf{$g_k$} & \textbf{$A_{ki}$ (s$^{-1}$)} \\
        \midrule
        \multicolumn{7}{l}{\textbf{Al I (Z = 6.04)}} \\
        \quad - & 396.152 & 0.014 & 3.143 & 4 & 2 & $9.85 \times 10^7$ \\
        \quad - & 309.271 & 0.014 & 4.022 & 4 & 6 & $7.29 \times 10^7$ \\
        \quad - & 394.401 & 0.000 & 3.143 & 2 & 2 & $4.99 \times 10^7$ \\
        \quad - & 308.215 & 0.000 & 4.021 & 2 & 4 & $5.87 \times 10^7$ \\
        \quad - & 237.312 & 0.014 & 5.237 & 4 & 6 & $9.07 \times 10^7$ \\
        \midrule
        \multicolumn{7}{l}{\textbf{Al II (Z = 1.01)}} \\
        \quad - & 281.618 & 7.421 & 11.822 & 3 & 1 & $3.57 \times 10^8$ \\
        \quad - & 266.916 & 0.000 & 4.644 & 1 & 3 & $3.28 \times 10^3$ \\
        \quad - & 704.208 & 11.317 & 13.077 & 3 & 5 & $5.78 \times 10^7$ \\
        \quad - & 705.671 & 11.317 & 13.073 & 3 & 3 & $5.74 \times 10^7$ \\
        \quad - & 466.305 & 10.598 & 13.256 & 5 & 3 & $5.81 \times 10^7$ \\
        \midrule
        \multicolumn{7}{l}{\textbf{Ca I (Z = 2.17)}} \\
        \quad - & 422.673 & 0.000 & 2.933 & 1 & 3 & $2.18 \times 10^8$ \\
        \quad - & 445.478 & 1.899 & 4.681 & 5 & 7 & $8.70 \times 10^7$ \\
        \quad - & 430.253 & 1.899 & 4.780 & 5 & 5 & $1.36 \times 10^8$ \\
        \quad - & 616.217 & 1.899 & 3.910 & 5 & 3 & $4.77 \times 10^7$ \\
        \quad - & 443.496 & 1.886 & 4.681 & 3 & 5 & $6.70 \times 10^7$ \\
        \midrule
        \multicolumn{7}{l}{\textbf{Ca II (Z = 2.92)}} \\
        \quad - & 393.366 & 0.000 & 3.151 & 2 & 4 & $1.47 \times 10^8$ \\
        \quad - & 396.847 & 0.000 & 3.123 & 2 & 2 & $1.40 \times 10^8$ \\
        \quad - & 854.209 & 1.700 & 3.151 & 6 & 4 & $9.90 \times 10^6$ \\
        \quad - & 866.214 & 1.692 & 3.123 & 4 & 2 & $1.06 \times 10^7$ \\
        \quad - & 317.933 & 3.151 & 7.050 & 4 & 6 & $3.60 \times 10^8$ \\
        \bottomrule
    \end{tabular}
\end{table}

Langkah terakhir adalah menggabungkan informasi makroskopis (fraksi Saha) dan mikroskopis (parameter transisi dan Z) untuk menghitung intensitas relatif final dari setiap garis, seperti yang dirinci pada Tabel~\ref{tab:intensity_calculation}. Kolom $I_{\text{Boltzmann}}$ adalah intensitas intrinsik transisi yang dihitung dari distribusi Boltzmann, yang kemudian diskalakan oleh Faktor Saha. Kolom terakhir menormalisasi semua intensitas final sehingga garis terkuat memiliki nilai 1.000.

\begin{table}[H]
    \centering
    \caption{Rincian perhitungan intensitas relatif untuk garis-garis utama Al dan Ca.}
    \label{tab:intensity_calculation}
    \small
    \begin{tabular}{lccccc}
        \toprule
        \textbf{Elemen} & \textbf{$\lambda$ (nm)} & \textbf{$I_{\text{Boltzmann}}$ (rel.)} & \textbf{Faktor Saha} & \textbf{$I_{\text{Final}}$ (rel.)} & \textbf{$I_{\text{Final}}$ (Normalized)} \\
        \midrule
        Al I  & 396.152 & $3.42 \times 10^5$ & 0.146 & $4.99 \times 10^4$ & 0.029 \\
        Al I  & 309.271 & $2.12 \times 10^5$ & 0.146 & $3.09 \times 10^4$ & 0.018 \\
        Al I  & 394.401 & $1.73 \times 10^5$ & 0.146 & $2.53 \times 10^4$ & 0.015 \\
        Al I  & 308.215 & $1.14 \times 10^5$ & 0.146 & $1.66 \times 10^4$ & 0.010 \\
        Al I  & 237.312 & $4.53 \times 10^4$ & 0.146 & $6.60 \times 10^3$ & 0.004 \\
        \midrule
        Al II & 281.618 & $1.26 \times 10^1$ & 0.854 & $1.08 \times 10^1$ & 0.000 \\
        Al II & 266.916 & $1.16 \times 10^1$ & 0.854 & $9.89 \times 10^0$ & 0.000 \\
        Al II & 704.208 & $1.66 \times 10^0$ & 0.854 & $1.41 \times 10^0$ & 0.000 \\
        Al II & 705.671 & $9.91 \times 10^{-1}$ & 0.854 & $8.47 \times 10^{-1}$ & 0.000 \\
        Al II & 466.305 & $7.69 \times 10^{-1}$ & 0.854 & $6.57 \times 10^{-1}$ & 0.000 \\
        \midrule
        Ca I  & 422.673 & $4.27 \times 10^6$ & 0.170 & $7.28 \times 10^5$ & 0.421 \\
        Ca I  & 445.478 & $3.15 \times 10^5$ & 0.170 & $5.37 \times 10^4$ & 0.031 \\
        Ca I  & 430.253 & $3.05 \times 10^5$ & 0.170 & $5.19 \times 10^4$ & 0.030 \\
        Ca I  & 616.217 & $2.26 \times 10^5$ & 0.170 & $3.86 \times 10^4$ & 0.022 \\
        Ca I  & 443.496 & $1.73 \times 10^5$ & 0.170 & $2.95 \times 10^4$ & 0.017 \\
        \midrule
        Ca II & 393.366 & $2.09 \times 10^6$ & 0.830 & $1.73 \times 10^6$ & \textbf{1.000} \\
        Ca II & 396.847 & $1.03 \times 10^6$ & 0.830 & $8.58 \times 10^5$ & 0.496 \\
        Ca II & 854.209 & $1.41 \times 10^5$ & 0.830 & $1.17 \times 10^5$ & 0.067 \\
        Ca II & 866.214 & $7.83 \times 10^4$ & 0.830 & $6.50 \times 10^4$ & 0.038 \\
        Ca II & 317.933 & $2.68 \times 10^4$ & 0.830 & $2.23 \times 10^4$ & 0.013 \\
        \bottomrule
    \end{tabular}
\end{table}
\subsection{Analisis Fisis dan Konfirmasi Visual}

Analisis dari Tabel~\ref{tab:intensity_calculation} memberikan wawasan fisis yang penting. Untuk Aluminium, terjadi "persaingan" antara populasi dan eksitasi. Garis Al I memiliki intensitas Boltzmann yang tinggi karena energi eksitasinya ($E_k$) rendah, namun kontribusinya dilemahkan oleh fraksi Saha yang kecil (hanya 14.6\%). Sebaliknya, garis Al II mendapat keuntungan dari fraksi Saha yang dominan (85.4\%), namun intensitas Boltzmann-nya sangat rendah akibat energi eksitasi yang sangat tinggi. Akibatnya, tidak ada garis Al yang menjadi sangat dominan pada kondisi ini.

Untuk Kalsium, ceritanya berbeda. Garis Ca II pada 393.366 nm dan 396.847 nm (garis H dan K) mendapat keuntungan ganda: intensitas Boltzmann yang tinggi (karena $E_k$ relatif rendah untuk ion) dan dikalikan dengan faktor Saha yang sangat besar (83.0\%). Kombinasi inilah yang menjadikan garis Ca II pada 393.366 nm sebagai garis emisi terkuat secara absolut. Menariknya, garis Ca I pada 422.673 nm juga cukup kuat (42.1\% dari maksimum), menunjukkan betapa besarnya intensitas Boltzmann intrinsiknya meskipun berasal dari spesies minoritas. Perhitungan ini secara langsung memprediksi fitur-fitur spektral utama yang akan diamati secara visual pada bagian selanjutnya.

\begin{figure}[H]
    \centering
    \includegraphics[width=0.9\textwidth]{images/Al-Ca-delta-electron-n_T8000K_ne1e17_wl200-900nm.png}
    \caption{Spektrum emisi komposit yang dihasilkan untuk campuran Al dan Ca pada T=8000 K dan $n_e=1.0 \times 10^{17}$ cm$^{-3}$. Plot ini secara visual mengkonfirmasi hasil perhitungan, menunjukkan dominasi kuat dari garis Ca II (393.366 nm dan 396.847 nm) dan garis Ca I (422.673 nm), sementara garis Al I (sekitar 394-396 nm) terlihat jauh lebih lemah.}
    \label{fig:al_ca_spectrum}
\end{figure}

Gambar~\ref{fig:al_ca_spectrum} menyajikan visualisasi akhir dari spektrum komposit yang dihasilkan dari semua perhitungan ini. Plot tersebut berfungsi sebagai konfirmasi visual dari analisis numerik: garis-garis Ca II pada 393.366 nm dan 396.847 nm secara jelas mendominasi spektrum, diikuti oleh garis Ca I pada 422.673 nm. Garis-garis dari Al I, meskipun ada, memiliki intensitas yang jauh lebih rendah, sesuai dengan prediksi. Studi kasus ini menunjukkan bagaimana interaksi kompleks antara parameter atomik dan kondisi plasma secara langsung membentuk pola emisi spektral yang diamati, dan memvalidasi bahwa proses simulasi secara akurat menangkap dinamika fundamental ini.

\section{Realisme Fisis pada Dataset Spektral Sintetis}
\label{sec:realisme_fisis}
Keberhasilan pendekatan yang diusulkan sangat bergantung pada kualitas dataset sintetis. Oleh karena itu, langkah pertama yang paling fundamental adalah memvalidasi dataset untuk memastikan ia merepresentasikan spektrum secara akurat dan mampu mereplikasi kompleksitas fisis yang ditemukan pada spektrum eksperimental.

\subsection{Kesesuaian Spektral dengan Standar Referensi (NIST)}
Pada dasarnya, validasi yang paling mendasar dilakukan dengan membandingkan hasil yang diperoleh terhadap data dari basis data \textit{NIST Atomic Spectra Database} (ASD), yang merupakan standar referensi dalam bidang spektroskopi. Perbandingan visual pada Gambar~\ref{fig:nist_comparison} mengindikasikan adanya kesesuaian yang sangat tinggi antara spektrum yang dihasilkan dengan spektrum referensi dari NIST, baik dari segi posisi puncak maupun intensitas relatifnya. Validasi ini memberikan keyakinan bahwa implementasi persamaan Saha-Boltzmann dan profil Gaussian dalam metode ini telah akurat dan dapat diandalkan secara fisis.

\begin{figure}[H]
    \centering
    \subfloat[Spektrum yang Dihasilkan (212-220 nm)]{\includegraphics[width=0.48\textwidth]{images/Al-Ca-Fe-Si_T8000K_ne1e17_wl212-220nm_dk.png}}
    \hfill
    \subfloat[Referensi NIST (212-220 nm)]{\includegraphics[width=0.48\textwidth]{images/Al-Ca-Fe-Si_T8000K_ne1e17-NIST-1.png}}
    \\
    \subfloat[Spektrum yang Dihasilkan (390-400 nm)]{\includegraphics[width=0.48\textwidth]{images/Al-Ca-Fe-Si_T8000K_ne1e17_wl390-400nm_dk.png}}
    \hfill
    \subfloat[Referensi NIST (390-400 nm)]{\includegraphics[width=0.48\textwidth]{images/Al-Ca-Fe-Si_T8000K_ne1e17-NIST.png}}
    \caption{Perbandingan spektrum emisi yang dihasilkan dengan referensi dari NIST ASD untuk campuran Al, Ca, Fe, Si pada $T = 8000$ K dan $n_e = 1.0 \times 10^{17}$ cm$^{-3}$.}
    \label{fig:nist_comparison}
\end{figure}

\subsection{Pemodelan Fenomena Fisika Plasma Kompleks}
Setelah akurasi dasar divalidasi, tingkat analisis selanjutnya adalah menunjukkan bahwa metode yang digunakan mampu mereplikasi fenomena fisika plasma yang dinamis dan kompleks.

\subsubsection{Interferensi Spektral Multi-Elemen}
Salah satu tantangan fisis utama dalam analisis LIBS adalah adanya interferensi spektral, di mana garis emisi dari elemen yang berbeda tumpang-tindih. Gambar~\ref{fig:interference} mendemonstrasikan kemampuan metode ini dalam memodelkan fenomena tersebut, sehingga merepresentasikan skenario realistis yang memungkinkan model DL untuk mempelajari cara memisahkan sinyal-sinyal yang saling mengganggu.

\begin{figure}[H]
    \centering
    \subfloat[Spektrum yang Dihasilkan (tanpa dekonvolusi)]{\includegraphics[width=0.48\textwidth]{images/Al-Ca-Fe-Si_T8000K_ne1e17_wl212-220nm.png}}
    \hfill
    \subfloat[Spektrum yang Dihasilkan (dengan dekonvolusi)]{\includegraphics[width=0.48\textwidth]{images/Al-Ca-Fe-Si_T8000K_ne1e17_wl212-220nm_dk.png}}
    \caption{Perbandingan spektrum emisi yang dihasilkan dengan dan tanpa dekonvolusi untuk campuran Al, Ca, Fe, Si.}
    \label{fig:interference}
\end{figure}

Untuk lebih mengilustrasikan kompleksitas ini, Gambar~\ref{fig:interference} menyajikan visualisasi spektrum dari data referensi NIST itu sendiri, yang menunjukkan bagaimana kontribusi dari berbagai tingkat ionisasi (seperti Al I, Al II, Ca I, dan Ca II) secara bersama-sama membentuk spektrum komposit akhir. Penggunaan skala logaritmik pada sumbu intensitas sangat penting di sini, karena memungkinkan visualisasi garis-garis emisi yang jauh lebih lemah yang mungkin tertutupi dalam skala linear, memberikan gambaran yang lebih lengkap tentang kekayaan informasi dalam spektrum.


\subsubsection{Dinamika Kesetimbangan Ionisasi Saha}
Aspek fundamental lain dari fisika plasma adalah kesetimbangan ionisasi yang diatur oleh persamaan Saha. Untuk menganalisis dinamika ini secara lebih kuantitatif, Tabel~\ref{tab:analisis_parametrik_suhu} menyajikan hasil analisis parametrik yang menunjukkan bagaimana intensitas relatif dari garis-garis spektral utama untuk beberapa elemen representatif berubah seiring dengan meningkatnya suhu.

\begin{table}[H]
    \centering
    \caption{Hasil Analisis Parametrik: Intensitas Relatif Ternormalisasi Berdasarkan Suhu. Intensitas dinormalisasi ke nilai maksimum pada setiap suhu untuk setiap elemen.}
    \label{tab:analisis_parametrik_suhu}
    \small
    \begin{tabular}{llrrrr}
        \toprule
        \multicolumn{2}{l}{\textbf{Elemen Garis Spektral}} & \multicolumn{4}{c}{\textbf{Suhu (K)}} \\
        \cmidrule(lr){3-6}
        & & \textbf{6000} & \textbf{8000} & \textbf{10000} & \textbf{12000} \\
        \midrule
        \multirow{6}{*}{Al} & Al I 309.27 nm   & 0.000 & 0.620 & 0.801 & 0.949 \\
                            & Al I 394.40 nm   & 0.507 & 0.000 & 0.000 & 0.000 \\
                            & Al I 396.15 nm   & 1.000 & 1.000 & 1.000 & 1.000 \\
                            & Al II 266.92 nm  & 0.000 & 0.000 & 0.000 & 0.000 \\
                            & Al II 281.62 nm  & 0.000 & 0.000 & 0.022 & 0.513 \\
                            & Al II 704.21 nm  & 0.000 & 0.000 & 0.004 & 0.123 \\
        \midrule
        \multirow{6}{*}{Ca} & Ca I 422.67 nm   & 1.000 & 0.421 & 0.037 & 0.007 \\
                            & Ca I 430.25 nm   & 0.000 & 0.000 & 0.000 & 0.001 \\
                            & Ca I 445.48 nm   & 0.000 & 0.031 & 0.005 & 0.000 \\
                            & Ca I 616.22 nm   & 0.033 & 0.000 & 0.000 & 0.000 \\
                            & Ca II 393.37 nm  & 0.056 & 1.000 & 1.000 & 1.000 \\
                            & Ca II 396.85 nm  & 0.028 & 0.496 & 0.492 & 0.489 \\
        \midrule
        \multirow{5}{*}{Fe} & Fe I 248.33 nm   & 1.000 & 1.000 & 0.395 & 0.056 \\
                            & Fe I 248.81 nm   & 0.000 & 0.674 & 0.270 & 0.039 \\
                            & Fe I 358.12 nm   & 0.920 & 0.000 & 0.000 & 0.000 \\
                            & Fe II 238.20 nm  & 0.002 & 0.143 & 0.965 & 1.000 \\
                            & Fe II 259.94 nm  & 0.002 & 0.169 & 1.000 & 0.953 \\
        \midrule
        \multirow{4}{*}{Mg} & Mg I 285.21 nm   & 1.000 & 1.000 & 0.318 & 0.042 \\
                            & Mg I 383.83 nm   & 0.035 & 0.075 & 0.038 & 0.007 \\
                            & Mg II 279.55 nm  & 0.003 & 0.194 & 1.000 & 1.000 \\
                            & Mg II 280.27 nm  & 0.001 & 0.098 & 0.501 & 0.500 \\
        \midrule
        \multirow{4}{*}{Si} & Si I 251.61 nm   & 1.000 & 1.000 & 0.941 & 0.155 \\
                            & Si I 288.16 nm   & 0.604 & 0.643 & 0.628 & 0.106 \\
                            & Si II 385.60 nm  & 0.001 & 0.057 & 0.753 & 0.753 \\
                            & Si II 634.71 nm  & 0.001 & 0.076 & 1.000 & 1.000 \\
        \bottomrule
    \end{tabular}
\end{table}

Data pada tabel ini secara jelas mengilustrasikan titik transisi di mana dominasi emisi bergeser dari spesies netral (I) ke spesies terionisasi (II). Misalnya, untuk Ca pada 6000 K, garis Ca I 422.67 nm adalah yang terkuat (nilai 1.000). Namun, pada 8000 K, intensitasnya turun drastis menjadi 0.421 sementara garis Ca II 393.37 nm mengambil alih sebagai yang terkuat. Pola pergeseran dominasi ini, dari garis atom netral ke garis ion seiring meningkatnya suhu, terlihat konsisten di semua elemen. Analisis kuantitatif ini, yang didukung oleh visualisasi pada Gambar~\ref{fig:saha_ionization_al} dan \ref{fig:saha_ionization_ca}, memperkuat validasi dataset dengan menunjukkan bahwa ia secara akurat memodelkan perilaku fisis fundamental yang menjadi dasar identifikasi unsur.

\begin{figure}[H]
    \centering
    \subfloat[Spektrum Al pada 6000 K (Dominan Al I)]{\includegraphics[width=0.8\textwidth]{images/al_spectrum_6000K.png}} \\
    \subfloat[Spektrum Al pada 30000 K (Dominan Al II)]{\includegraphics[width=0.8\textwidth]{images/al_spectrum_30000K.png}}
    \caption{Ilustrasi efek kesetimbangan ionisasi Saha pada spektrum Al murni.}
    \label{fig:saha_ionization_al}
\end{figure}

\begin{figure}[H]
    \centering
    \subfloat[Spektrum Ca pada 6000 K (Dominan Ca I)]{\includegraphics[width=0.8\textwidth]{images/ca_spectrum_6000K.png}} \\
    \subfloat[Spektrum Ca pada 8000 K (Dominan Ca II)]{\includegraphics[width=0.8\textwidth]{images/ca_spectrum_8000K.png}}
    \caption{Ilustrasi efek kesetimbangan ionisasi Saha pada spektrum Ca murni.}
    \label{fig:saha_ionization_ca}
\end{figure}

\section{Analisis Kemampuan Model dalam Menginterpretasi Fisika Spektral}
\label{sec:interpretasi_fisis_model}
Setelah fondasi data divalidasi, analisis beralih pada kemampuan model \textit{Informer} dalam mempelajari dan menginterpretasikan fisika dari data spektral. Pembahasan disajikan dari bukti yang paling sederhana hingga analisis kuantitatif yang paling kompleks.

\subsection{Stabilitas Proses Pembelajaran terhadap Data Fisis}
Langkah pertama dalam menganalisis kinerja model adalah memastikan proses pelatihannya berlangsung secara stabil. Gambar~\ref{fig:training_loss} menampilkan kurva pembelajaran yang mengindikasikan bahwa baik \textit{loss} pelatihan maupun validasi menurun secara konsisten dan konvergen. Hal ini menunjukkan bahwa proses pelatihan model bersifat stabil dan berhasil menangkap struktur fisis yang terkandung di dalam data tanpa mengalami \textit{overfitting} yang signifikan.

\begin{figure}[H]
    \centering
    \includegraphics[width=0.8\textwidth]{images/training_history_plot.png}
    \caption{Grafik kurva pembelajaran model \textit{Informer} (3-layer) selama 70 epoch. Penurunan yang stabil pada kurva loss dan peningkatan F1-Score mengkonfirmasi proses pembelajaran yang efektif.}
    \label{fig:training_loss}
\end{figure}

\subsection{Demonstrasi Dekonvolusi Spektral Kualitatif oleh Model}
Setelah proses pembelajaran terbukti stabil, selanjutnya disajikan hasil kerja model secara kualitatif. Gambar~\ref{fig:prediction_sample_10} dan \ref{fig:prediction_grup_1} merepresentasikan hasil prediksi model pada data uji sintetis dan data eksperimental. Visualisasi ini secara intuitif menunjukkan kemampuan model untuk melakukan dekomposisi spektral secara kualitatif---yaitu memisahkan spektrum komposit yang kompleks dan memberikan label pada puncak-puncak spektral sesuai elemen penyebabnya. Kemampuan ini menjadi bukti visual yang kuat dari efektivitas pendekatan yang diusulkan.

\begin{figure}[H]
    \centering
    \subfloat[200nm - 375nm]{\includegraphics[width=0.48\textwidth]{images/prediction_plot_sample_10_range_1.png}}
    \hfill
    \subfloat[375nm - 550nm]{\includegraphics[width=0.48\textwidth]{images/prediction_plot_sample_10_range_2.png}}
    \\
    \subfloat[550nm - 725nm]{\includegraphics[width=0.48\textwidth]{images/prediction_plot_sample_10_range_3.png}}
    \hfill
    \subfloat[725nm - 900nm]{\includegraphics[width=0.48\textwidth]{images/prediction_plot_sample_10_range_4.png}}
    \caption{Visualisasi prediksi model 2-layer pada Sampel 10 dari dataset uji.}
    \label{fig:prediction_sample_10}
\end{figure}

\begin{figure}[H]
    \centering
    \subfloat[200nm - 375nm]{\includegraphics[width=0.48\textwidth]{images/prediction_plot_Grup-1_range_1.png}}
    \hfill
    \subfloat[375nm - 550nm]{\includegraphics[width=0.48\textwidth]{images/prediction_plot_Grup-1_range_2.png}}
    \\
    \subfloat[550nm - 725nm]{\includegraphics[width=0.48\textwidth]{images/prediction_plot_Grup-1_range_3.png}}
    \hfill
    \subfloat[725nm - 900nm]{\includegraphics[width=0.48\textwidth]{images/prediction_plot_Grup-1_range_4.png}}
    \caption{Visualisasi prediksi model 2-layer pada data lapangan dari sampel cangkang telur.}
    \label{fig:prediction_grup_1}
\end{figure}

\subsection{Interpretasi Fisis Kinerja Prediktif: Sensitivitas vs. Selektivitas}
Setelah bukti kualitatif disajikan, analisis paling kompleks adalah evaluasi kuantitatif beserta interpretasi fisisnya. Tabel~\ref{tab:perbandingan_lengkap} menyajikan perbandingan kinerja detail antara arsitektur model 2-layer dan 3-layer.
\begin{table}[H]
    \centering
    \begin{threeparttable}
    \small
    \caption{Perbandingan Kinerja Detail Model Informer 2-Layer dan 3-Layer pada Dataset Uji.}
    \label{tab:perbandingan_lengkap}
    \begin{tabular}{l rrr rrr r}
        \toprule
        & \multicolumn{3}{c}{\textbf{Model 2-Layer (Optimal)}} & \multicolumn{3}{c}{\textbf{Model 3-Layer}} & \\
        \cmidrule(lr){2-4} \cmidrule(lr){5-7}
        \textbf{Elemen} & \textbf{Prec} & \textbf{Rec} & \textbf{F1} & \textbf{Prec} & \textbf{Rec} & \textbf{F1} & \textbf{\textit{Support}} \\
        \midrule
        Al & \textbf{0.5217} & \textbf{0.9225} & \textbf{0.6665} & 0.4302 & 0.3970 & 0.4130 & 2174388 \\
        Ar & \textbf{0.5007} & \textbf{0.5605} & \textbf{0.5289} & 0.3570 & 0.3580 & 0.3575 & 3016095 \\
        C  & \textbf{0.4391} & \textbf{0.5466} & \textbf{0.4870} & 0.3469 & 0.3835 & 0.3643 & 1729930 \\
        Ca & \textbf{0.6758} & \textbf{1.0000} & \textbf{0.8065} & 0.5066 & 0.9992 & 0.6724 & 1127079 \\
        Cl & \textbf{0.5291} & 0.9996 & \textbf{0.6920} & 0.3831 & \textbf{0.9999} & 0.5540 & 189973 \\
        Co & \textbf{0.5516} & 0.2016 & 0.2953 & 0.5065 & \textbf{0.3201} & \textbf{0.3923} & 4381508 \\
        Cr & \textbf{0.7245} & \textbf{0.9972} & \textbf{0.8392} & 0.5106 & 0.8987 & 0.6512 & 2983956 \\
        Fe & 0.5783 & 0.1038 & 0.1761 & \textbf{0.6287} & \textbf{0.2744} & \textbf{0.3821} & 6379378 \\
        Mg & \textbf{0.6252} & \textbf{0.9813} & \textbf{0.7638} & 0.4541 & 0.7121 & 0.5545 & 2159572 \\
        Mn & \textbf{0.7500} & \textbf{0.8874} & \textbf{0.8129} & 0.6045 & 0.6496 & 0.6262 & 4939082 \\
        N  & \textbf{0.5104} & \textbf{0.6495} & \textbf{0.5716} & 0.3414 & 0.4870 & 0.4014 & 2089132 \\
        Na & \textbf{0.6950} & \textbf{0.9978} & \textbf{0.8193} & 0.5239 & 0.9041 & 0.6634 & 2442502 \\
        Ni & \textbf{0.7226} & \textbf{0.9758} & \textbf{0.8303} & 0.5971 & 0.8565 & 0.7037 & 3582470 \\
        O  & \textbf{0.3830} & \textbf{0.7988} & \textbf{0.5178} & 0.2614 & 0.5882 & 0.3619 & 979582 \\
        S  & \textbf{0.5721} & \textbf{0.9932} & \textbf{0.7260} & 0.3772 & 0.9257 & 0.5360 & 1169165 \\
        Si & \textbf{0.7244} & \textbf{0.9986} & \textbf{0.8397} & 0.4183 & 0.9770 & 0.5858 & 1565405 \\
        Ti & \textbf{0.5876} & \textbf{0.5046} & \textbf{0.5430} & 0.5621 & 0.4885 & 0.5227 & 5288926 \\
        background & 0.5981 & \textbf{0.8008} & \textbf{0.6848} & \textbf{0.5998} & 0.6314 & 0.6152 & 8693531 \\
        \midrule
        \textbf{\textit{micro avg}}    & \textbf{0.6142} & \textbf{0.6815} & \textbf{0.6461} & 0.5011 & 0.5840 & 0.5394 & 54891674 \\
        \textbf{\textit{macro avg}}    & \textbf{0.5938} & \textbf{0.7733} & \textbf{0.6445} & 0.4672 & 0.6584 & 0.5199 & 54891674 \\
        \textbf{\textit{weighted avg}} & \textbf{0.6090} & \textbf{0.6815} & \textbf{0.6044} & 0.5219 & 0.5840 & 0.5256 & 54891674 \\
        \bottomrule
    \end{tabular}
    \begin{tablenotes}
      \small
      \item \textit{Catatan:} Nilai yang dicetak \textbf{tebal} menyoroti kinerja yang lebih unggul antara kedua model pada setiap metrik.
    \end{tablenotes}
    \end{threeparttable}
\end{table}

Dari tabel ini, arsitektur 2-layer secara jelas menunjukkan kinerja yang jauh lebih unggul pada metrik F1-Score (Macro) sebesar \textbf{0.6445}, yang mengindikasikan tercapainya keseimbangan optimal antara kompleksitas dan kemampuan generalisasi. Interpretasi fisis yang lebih mendalam dari pola metrik tetap relevan:
\begin{itemize}
    \item \textbf{Recall Tinggi (Sensitivitas Tinggi):} menunjukkan bahwa model sangat sensitif dalam mendeteksi pola emisi spektral setiap elemen.
    \item \textbf{Precision Rendah (Selektivitas Terbatas):} Di sisi lain, nilai \textit{precision} yang lebih rendah, terutama untuk elemen dengan sinyal kompleks seperti Co dan Fe, merepresentasikan kesulitan fisis nyata dalam membedakan sinyal asli dari derau atau interferensi spektral yang tumpang-tindih.
\end{itemize}

\section{Diskusi Implikasi Fisis dan Metodologis}
\label{sec:diskusi_implikasi}
Bagian terakhir ini membahas makna yang lebih luas dari temuan yang telah disajikan, dengan fokus pada implikasinya bagi fisika eksperimental dan metodologi analisis data spektral.

\subsection{Implikasi Fisis: Menuju Paradigma Pengukuran Bebas Kalibrasi}
Keberhasilan model yang dilatih murni pada data sintetis dalam mencapai kinerja prediksi yang tinggi merupakan bukti konsep yang kuat untuk pendekatan analisis LIBS bebas kalibrasi. Temuan ini mengindikasikan bahwa model berhasil mempelajari hubungan fundamental antara parameter fisis plasma dan pola emisi spektral yang dihasilkannya. Hal ini berpotensi mengatasi ketergantungan pada sampel standar bersertifikat yang mahal dan memakan waktu.

\subsection{Refleksi Metodologis: Peran Teknik Komputasi dalam Menerjemahkan Fisika}
Pencapaian hasil ini juga didukung oleh teknik pelatihan lanjutan yang membantu model menginterpretasikan fisika spektral secara efektif:
\begin{itemize}
    \item \textbf{Focal Loss:} Teknik ini memungkinkan model untuk lebih fokus pada puncak emisi yang lemah atau tumpang-tindih, yang secara fisis sulit dideteksi, sehingga berkontribusi pada peningkatan \textit{recall}.
    \item \textbf{Regularisasi:} Implementasi regularisasi mencegah model mengalami \textit{overfitting} atau menemukan korelasi spurius yang palsu, yang dibuktikan oleh kurva pembelajaran yang baik (Gambar~\ref{fig:training_loss}).
    \item \textbf{\textit{Automatic Mixed Precision} (AMP):} Pada dasarnya, teknik ini merupakan kebutuhan praktis yang memungkinkan dilakukannya eksperimen komputasi skala besar yang diperlukan untuk memvalidasi pendekatan ini secara kokoh.
\end{itemize}

\subsection{Keterbatasan dari Perspektif Model Fisika}
Meskipun menjanjikan, pendekatan ini memiliki keterbatasan yang perlu diakui dari sudut pandang fisika:
\begin{enumerate}
    \item \textbf{Penyederhanaan Model Fisika:} Dataset sintetis didasarkan pada model fisika yang disederhanakan, seperti asumsi LTE dan pengabaian efek matriks non-linear yang kompleks, yang mungkin tidak sepenuhnya merepresentasikan kondisi eksperimental riil.
    \item \textbf{Keterbatasan Pengetahuan Model:} Model saat ini hanya dilatih untuk mengenali fisika dari 18 elemen yang dilatihkan, dan kemampuannya untuk mengidentifikasi elemen di luar set ini belum teruji.
    \item \textbf{Validitas pada Kondisi Non-LTE:} Kinerja model pada spektrum yang berasal dari plasma non-LTE, yang dapat terjadi pada fase awal evolusi plasma LIBS, belum dievaluasi.
\end{enumerate}